\documentclass{ctexart}

\usepackage{palatino}
\usepackage{tikz}
\usetikzlibrary{shapes.geometric, arrows}
\usepackage{bm}
\usepackage{geometry}
\usepackage{titlesec} %自u格式的宏包
\usepackage{picinpar,graphicx}
\usepackage{zhnumber}
\usepackage{float}
\usepackage{caption,subcaption}
\usepackage{amsmath}
\usepackage{amssymb}
\usepackage{fancyhdr}%页眉页脚
\usepackage{multirow}
\usepackage{upgreek}
\usepackage{paralist}
\usepackage{diagbox,array}
\usepackage{cite}
\usepackage{minted}
\usepackage{listings}
\usepackage{xcolor}
\usepackage{hyperref}
\usepackage{tabularx}
\usepackage{booktabs}
\usepackage{fancyhdr}
\usepackage{verbatim}
\usepackage{supertabular}
\usepackage{subfiles}
\usepackage{mathtools}


\lstset{
	basicstyle=\small,%
	escapeinside=``,%
	keywordstyle=\color{red} \bfseries,% \underbar,%
	identifierstyle={},%
	commentstyle=\color{blue},%
	stringstyle=\ttfamily,%
	%labelstyle=\tiny,%
	extendedchars=false,%
	linewidth=\textwidth,%
	numbers=left,%
	numberstyle=\tiny \color{blue},%
	frame=trbl%
}

\geometry{a4paper,left=0.75in,right=0.75in,top=1in,bottom=1in} %页边距


\counterwithin{figure}{section}
\counterwithin{table}{section}
\counterwithin{equation}{section}

\renewcommand{\thefigure}{ \arabic{section}-\arabic{figure}}
\renewcommand{\thetable}{ \arabic{section}-\arabic{table}}
\renewcommand{\theequation}{ \arabic{section}-\arabic{equation}}
\renewcommand \thesubsection {(\zhnum{section})}
\renewcommand{\caption}[1]{\caption[\songti\ziju{-5}]{#1}}
%\\totalnumber{10}
%\renewcommand{\thesubfigure}{\arabic{section}-\arabic{figure}-\arabic{subfigure}}

\titleformat{\section}[block]{\zihao{4}\bfseries}{\arabic{section}.}{1em}{}[]
\titleformat{\subsection}[block]{\zihao{4}\bfseries}{\quad\arabic{section}.\arabic{subsection}.}{1em}{}[]
\titleformat{\subsubsection}[block]{\zihao{4}\bfseries}{\quad\quad\arabic{section}.\arabic{subsection}.\arabic{subsubsection}.}{1em}{}[]
\titleformat{\paragraph}[block]{\zihao{-4}\bfseries}{\quad\quad\arabic{section}.\arabic{subsection}.\arabic{subsubsection}.\arabic{paragraph}.}{0em}{}[]

\setlength{\baselineskip}{22pt}



\zihao{-4}
%\setlength{\parskip}{0.1pt}
\bibliographystyle{IEEEtran}

\pagestyle{fancy}%清除原页眉页脚样式
\fancyhf{} 
\fancyhead[C]{2024年测控仪器电路三级项目报告}

\rfoot[C]{\centering\thepage}

\begin{document}
	
%%%%%%%%%%%%%第一题:供电电路%%%%%%%%%%%%%			完成
	\section{供电电路}
		\subfile{./各题论文编译文件/1.第一题}
		
%%%%%%%%%%%%%第二题:电压跟随器%%%%%%%%%%%%%			完成
	\newpage
	\section{电压跟随器}
		\subfile{./各题论文编译文件/2.第二题}
		
%%%%%%%%%%%%%第三题:同相放大电路%%%%%%%%%%%%%		完成
	\newpage
	\section{同相放大电路}
		\subfile{./各题论文编译文件/3.第三题}

%%%%%%%%%%%%%第四题:反相放大电路%%%%%%%%%%%%%		完成
	\newpage
	\section{反相放大电路}
		\subfile{./各题论文编译文件/4.第四题}
		
%%%%%%%%%%%%%第五题:反相放大电路%%%%%%%%%%%%%		完成
	\newpage
	\section{V/I转换器、I/V转换器和电流放大器}
		\subfile{./各题论文编译文件/5.第五题}

%%%%%%%%%%%%%第六题:反相放大电路%%%%%%%%%%%%%		缺少现象评价
	\newpage
	\section{加法器电路}
		\subfile{./各题论文编译文件/6.第六题}		
			
%%%%%%%%%%%%%第七题:反相放大电路%%%%%%%%%%%%%		完成			
	\newpage
	\section{一阶RC无源低通、高通滤波器}
		\subfile{./各题论文编译文件/7.第七题}		

%%%%%%%%%%%%%第八题:同相一阶有源高通、低通滤波器%%%%%%%%%%%%%		完成
	\newpage
	\section{同相一阶有源高通、低通滤波器}
		\subfile{./各题论文编译文件/8.第八题}		

%%%%%%%%%%%%%第九题:反相一阶有源高通、低通滤波器%%%%%%%%%%%%%		完成
	\newpage
	\section{反相一阶有源高通、低通滤波器}
		\subfile{./各题论文编译文件/9.第九题}

%%%%%%%%%%%%%第十题:Sellen-Key结构二阶有源低通滤波器、二阶有源高通滤波器%%%%%%%%%%%%%	完成
	\newpage
	\section{Sallen-Key结构二阶有源低通滤波器、二阶有源高通滤波器}
		\subfile{./各题论文编译文件/10.第十题}

%%%%%%%%%%%%%第十一题:无源带通滤波器%%%%%%%%%%%%%	完成
	\newpage
	\section{无源带通滤波器}
		\subfile{./各题论文编译文件/11.第十一题}	

%%%%%%%%%%%%%第十二题:由一阶有源高通-低通滤波器构成的二阶有源带通滤波器%%%%%%%%%%%%%	完成			
	\newpage
	\section{由一阶有源高通-低通滤波器构成的二阶有源带通滤波器}
		\subfile{./各题论文编译文件/12.第十二题}	

%%%%%%%%%%%%%第十三题:Sellen-Key结构二阶有源带通滤波器%%%%%%%%%%%%%		完成
	\newpage
	\section{Sallen-Key结构二阶有源带通滤波器}
		\subfile{./各题论文编译文件/13.第十三题}
		
%%%%%%%%%%%%%第十四题:Sellen-Key结构二阶有源带阻滤波器%%%%%%%%%%%%%		完成
	\newpage
	\section{Sallen-Key结构二阶有源带阻滤波器}
		\subfile{./各题论文编译文件/14.第十四题}	

%%%%%%%%%%%%%第十五题:MFB结构二阶有源带通滤波器%%%%%%%%%%%%%		完成
	\newpage
	\section{MFB结构二阶有源带通滤波器}
		\subfile{./各题论文编译文件/15.第十五题}

%%%%%%%%%%%%%第十六题:MFB结构二阶有源低通、高通滤波器%%%%%%%%%%%%%	
	\newpage
	\section{MFB结构二阶有源低通、高通滤波器}
		\subfile{./各题论文编译文件/16.第十六题}

%%%%%%%%%%%%%第十七题:三阶低通滤波器%%%%%%%%%%%%%		完成
	\newpage
	\section{三阶低通滤波器}
		\subfile{./各题论文编译文件/17.第十七题}	

%%%%%%%%%%%%%第十八题:四阶低通滤波器%%%%%%%%%%%%%				
	\newpage
	\section{四阶低通滤波器}
		\subfile{./各题论文编译文件/18.第十八题}	

%%%%%%%%%%%%%第十九题:调研市场中常用电测仪器设备%%%%%%%%%%%%%		完成
	\newpage
	\section{调研市场中常用电测仪器设备}
		\subfile{./各题论文编译文件/19.第十九题}	
\end{document}









